\section{Related work}

The availability of short-term measurements before primary outcomes motivates both surrogate analysis and methods for incomplete trial data. \citet{Prentice1989} formulates criteria for surrogate endpoints, while \citet{Fleming1996} emphasizes the substantive evidence required to connect surrogate responses to clinical effects. Short-term outcomes can also inform treatment-effect estimation while primary outcomes are still being collected, as in \citet{VanLancker2020}. Here, the surrogate enters the conditioning information for outcome arrival, and the target remains the population average treatment effect. Identification rests on randomized assignment, consistency, and outcome missingness at random conditional on baseline category, treatment, and surrogate, as specified in \cref{obj:ass:randomized-independence,obj:ass:surrogate-consistency,obj:ass:outcome-consistency,obj:ass:arrival-mar}.

The closest comparison for estimation rates is \citet[Assumption 2, Theorems 1--2, Section 3.2, and Appendix C.5]{ZengBalakrishnanHanKennedy2026}. Their observational experiment comprises independent observations of discrete covariates, treatment, and outcome, with treatment propensities bounded away from zero and one by a fixed positivity margin. Writing \leanref{sym:n}{\ensuremath{n}} for the sample size and \leanref{sym:d}{\ensuremath{d}} for the number of covariate categories, their plug-in analysis gives a squared-error upper bound of order \(n^{-1}+d^2/n^2\). Their minimax lower bound is of order \(n^{-1}+d^2/\{n^2\log^2 n\}\) in the stated regime \(d\lesssim n\log n\), with constants depending on positivity. These conclusions give a logarithmic comparison between the plug-in guarantee and the minimax benchmark. Appendix C.5 develops the moment-matching argument underlying that benchmark. For the corresponding observational regression contrast under fixed strict-interior positivity on occupied categories, \cref{obj:thm:zeng-fixed-positivity-polynomial-frontier} supplies a constructive polynomial-estimation guarantee through synthetic assignment. Its positivity margin, \leanref{sym:\epsilon}{\ensuremath{\epsilon}}, is held fixed in this comparison.

An earlier conference announcement by \citet[abstract 300411]{Kennedy2019DiscreteAbstract} considers an average causal effect with arbitrarily high-dimensional discrete confounders. It reports plug-in risk bounds, a minimax lower bound, and a variation with propensity-score positivity deteriorating with sample size, thereby establishing earlier consideration of the shrinking-positivity observational problem. The theorem-level comparison permitted by the available abstract is recorded in the limitations section.

The randomized analysis addresses a different observation mechanism. In the \leanref{S-1}{randomized fixed-horizon trial with a finite baseline alphabet}, treatment is assigned with known probability one half, and primary outcomes arrive with unknown probabilities conditional on the observed cells. The arrival floor \leanref{sym:q}{\ensuremath{q}} is imposed on occupied cells in \cref{obj:ass:occupied-cell-arrival}. Under the rare-arrival restriction in \cref{obj:ass:rare-arrival-slice}, \cref{obj:thm:matched-minimax-frontier} matches upper and lower squared-error bounds at the rate defined in \cref{obj:def:frontier-rate},
\[
\min\left\{1,\frac{1}{nq}
+\left[\frac{d}{nq\log(e+nq)}\right]^2\right\}.
\]
This comparison separates fixed treatment positivity in the observational experiment from outcome ascertainment in a balanced randomized experiment. The latter analysis quantifies the joint contribution of the effective sample size and the number of categories, while the synthetic-assignment construction connects the two experiments.

Surrogate-assisted efficiency provides a complementary perspective. \citet[Theorems 2.1--2.2 and 4.1--4.2]{KallusMao2025} characterize efficiency bounds with different patterns of surrogate and outcome observation, under treatment unconfoundedness, outcome missingness at random, and the overlap conditions appropriate to each observation pattern. Their labelled-sample asymptotics allow the labelled fraction to vanish while the expected number of labelled observations diverges, with fixed conditional distributions within the labelled and unlabelled populations. Attainment of the efficiency bound requires their nuisance convergence conditions: relevant products of regression and propensity or density-ratio errors must be negligible relative to the inverse square root of the expected labelled sample size. Treatment propensity estimates satisfy the stated boundedness condition, and estimated labelling propensities remain positive. Their moment requirements include uniformly bounded moments of order greater than two for the outcome residual, density ratio, and specified regression terms \citep[Assumption 7 and Appendix D.1, Assumption 9]{KallusMao2025}. These results explain how predictive surrogates improve asymptotic precision under those conditions. The finite-alphabet minimax analysis here characterizes uniform risk as category frequencies and the alphabet size vary over the stated model class.

Decaying outcome observation is also central to \citet[Theorems 2.2 and 4.1, equation (2.3)]{ZhangChakraborttyBradic2025}. Their triangular-array observational framework combines treatment overlap with positive, potentially vanishing labelling probabilities conditional on treatment and covariates. Theorem 2.2 establishes asymptotic normality when both nuisance models are correctly specified, the inverse-propensity effective sample size diverges, and the product-rate condition in equation (2.3) holds. Its conditions also include finite inverse-propensity expectations, bounded outcome variance, conditional residual variance bounded above and away from zero, and the stated weighted nuisance and tail requirements. Their bias-reduced estimator in Theorem 4.1 accommodates nuisance approximation error under its product restriction and alternative sparsity conditions, together with sub-Gaussian design and conditional noise assumptions, covariance lower bounds, propensity moment control, and bounded offset signals \citep[Assumptions 1--3 and 5--8]{ZhangChakraborttyBradic2025}. These conditions support asymptotic inference with estimated nuisance functions. Our randomized rare-arrival comparison instead measures worst-case squared error over a finite categorical model, and the interval analysis imposes uniform coverage over that same experiment.

The approximation methods have established antecedents in the estimation of discrete-distribution functionals. Polynomial estimators and moment-matched lower bounds are developed by \citet{JiaoVenkatHanWeissman2015}, \citet{Wu2016}, and \citet{WuYang2019}. Their pairing reflects the dual relationship between polynomial approximation and separation of distributions with matching moments. Classical approximation theory supplies the reciprocal-function approximation used in the converse argument \citep{Meinardus1967}; explicit reciprocal approximation formulas are also presented in \citet[Section 3]{https://doi.org/10.1007/s11075-026-02364-1}. The application here combines these tools with observed cell memberships and selectively observed outcomes. The roles of reciprocal approximation and moment-matched priors are stated in \cref{obj:lem:reciprocal-best-approximation,obj:lem:moment-matched-reciprocal-priors}, and the resulting comparison concerns the full observed experiment in \cref{obj:thm:full-experiment-lower}.

Bias correction for regression functionals supplies another comparison. Higher-order influence functions use higher-order U-statistics to obtain rate-optimal point and interval procedures for nonlinear functionals under smoothness and nuisance-estimation conditions \citep{RobinsLiTchetgenVanderVaart2008,Robins2009}. Augmented minimax linear estimation corrects a regression plug-in estimator by a minimax linear estimate of its error and obtains semiparametric efficiency under its function-class conditions \citep{Hirshberg2021}. Weak-overlap analysis instead thresholds doubly robust scores and derives asymptotic normality and Wald inference under propensity-tail, outcome-smoothness, and nuisance-rate restrictions \citep{Dorn2025}. The present model replaces smooth nuisance classes by arbitrary masses and heterogeneous arrival probabilities on a finite, growing alphabet. Its guarantees are nonasymptotic worst-case squared-error and honest-length comparisons, and its logarithmic improvement comes from polynomial correction of scarce arrived-cell counts rather than smooth regression approximation.

The interval criterion connects the analysis to uncertainty quantification through maximal expected length subject to uniform coverage. \citet{Donoho1994} relates statistical estimation to optimal recovery, \citet{Cai2004} studies expected-length comparisons and adaptation for confidence intervals, and \citet{ArmstrongKolesar2020} develops honest intervals that account for worst-case bias in nonparametric regression. For the present missing-outcome experiment, \cref{obj:def:interval-length-risk} specifies honesty for connected intervals, and \cref{obj:thm:connected-interval-frontier} gives the corresponding length comparison under rare outcome arrival. This places estimation error and honest interval length within a common finite-alphabet decision problem.

At complete outcome arrival, \citet{Sudijono2026} provides a useful comparison from randomized-experiment theory. They jointly optimize the assignment design and estimator for the sample average treatment effect in a finite population with binary potential outcomes, obtaining a sharp minimax analysis that includes a second-order risk expansion. Our complete-arrival result in \cref{obj:thm:unrestricted-complete-arrival-frontier} concerns a population average treatment effect under independent and identically distributed sampling and a fixed balanced assignment design. The common first-order sampling scale therefore arises in two distinct decision problems: joint design and estimation for a finite-population sample effect, and estimation of a population effect under a specified randomized sampling experiment.
